\vspace{0.5em}\noindent\textbf{FROM ZERO TO ZERO-DOWNTIME: DEVOPS \& CI/CD ON
AWS}\par

When starting with DevOps on AWS, CI/CD seemed like just automating
deployment scripts. But real Production systems taught that DevOps is
about \textbf{Reliability}, \textbf{Security}, and \textbf{Velocity}.

\vspace{0.5em}\noindent\textbf{3 Key Lessons}\par

\vspace{0.5em}\noindent\textbf{1. Infrastructure as Code (IaC): Say No to
ClickOps}\par

\begin{itemize}
\tightlist
\item
  Manual console clicks are fast initially but become a nightmare for
  environment replication and auditing.
\item
  \textbf{Use Terraform or AWS CDK}. Treat infrastructure as code. All
  changes go through Pull Request \(\rightarrow\) Plan \(\rightarrow\)
  Code Review.
\end{itemize}

\vspace{0.5em}\noindent\textbf{2. Shift-Left Security in CI/CD
Pipeline}\par

\begin{itemize}
\tightlist
\item
  Don't wait until Production to scan for vulnerabilities.
\item
  \textbf{Integrate security early:}

  \begin{itemize}
  \tightlist
  \item
    Scan IaC files with Checkov/TFLint
  \item
    Scan container images with Amazon ECR Image Scanning
  \item
    Manage secrets centrally with AWS Secrets Manager (never hardcode)
  \end{itemize}
\end{itemize}

\vspace{0.5em}\noindent\textbf{3. Blue/Green \& Canary Deployments for
Zero-Downtime}\par

\begin{itemize}
\tightlist
\item
  In-place deployments cause disruptions and 502 errors.
\item
  \textbf{Use AWS CodeDeploy:}

  \begin{itemize}
  \tightlist
  \item
    \textbf{Blue/Green}: Deploy parallel environment, test, then switch
    traffic
  \item
    \textbf{Canary}: Route 10\% traffic to new version, monitor
    CloudWatch metrics, then rollout 100\%
  \end{itemize}
\end{itemize}

\vspace{0.5em}\noindent\textbf{Detailed Pipeline
Architecture}\par

Below is the CI/CD workflow applied in practice:

\textbf{Step 1: Developer commits code to GitHub}

When developers push code to the repository, this is the starting point
of the entire pipeline.

\textbf{Step 2: AWS CodePipeline detects changes and triggers the
pipeline}

CodePipeline continuously monitors the repository. When a new commit is
detected, the pipeline is automatically triggered without manual
intervention.

\textbf{Step 3: AWS CodeBuild runs Unit Tests}

Code quality is verified through automated tests. If tests fail, the
pipeline stops immediately to prevent deploying broken code.

\textbf{Step 4: CodeBuild runs Security Scan}

Source code and dependencies are scanned to detect security
vulnerabilities. Early detection significantly reduces the cost of
fixing issues compared to finding them in Production.

\textbf{Step 5: Build container image and push to Amazon ECR}

A Docker image is built from the source code, then pushed to Amazon
Elastic Container Registry for storage and version management.

\textbf{Step 6: ECR automatically scans the image}

Once the image is pushed, ECR automatically scans it to detect CVEs in
the image layers. Scan results are recorded for evaluation before
deployment.

\textbf{Step 7: CodeDeploy (or Helm) deploys to Amazon EKS}

The image is deployed to the Kubernetes cluster. Depending on the setup,
either CodeDeploy (with native EKS integration) or Helm (for chart-based
management) can be used.

\textbf{Step 8: Apply Blue/Green or Canary strategy}

Instead of deploying directly to production, the pipeline uses one of
these strategies:

\begin{itemize}
\tightlist
\item
  Blue/Green: Creates a new Green version alongside the existing Blue
  version. After thorough testing on Green, traffic is switched from
  Blue to Green.
\item
  Canary: Routes a small percentage of traffic (e.g., 10\%) to the new
  version, monitors key metrics, then gradually rolls out to 100\%.
\end{itemize}

\textbf{Step 9: Monitor CloudWatch metrics and auto-rollback}

Throughout the deployment, CloudWatch tracks critical metrics like error
rates and response times. If anomalies are detected, the pipeline
automatically rollbacks to the previous version to minimize impact.

\vspace{0.5em}\noindent\textbf{Component Breakdown:}\par

\textbf{GitHub:} Source code repository with built-in CodePipeline
integration.

\textbf{AWS CodePipeline:} CI/CD orchestration service that coordinates
the entire workflow from build to deploy.

\textbf{AWS CodeBuild:} Serverless build service that runs tests,
security scans, and builds images.

\textbf{Amazon ECR:} Container image registry with automated CVE
scanning.

\textbf{AWS CodeDeploy / Helm:} Deployment tools with Zero-Downtime
strategies.

\textbf{Amazon EKS:} Kubernetes cluster where applications run.

\vspace{0.5em}\noindent\textbf{Results Achieved:}\par

\begin{itemize}
\tightlist
\item
  Deployment time reduced from 30 minutes to under 4 minutes
\item
  Early detection of security vulnerabilities at the build stage
\item
  Zero-downtime deployments through Blue/Green and Canary
\item
  Automatic rollback if issues are detected after deployment
\end{itemize}

\vspace{0.5em}\noindent\textbf{Key Takeaway}\par

Tools are only 20\%. The remaining 80\% is \textbf{Pipeline Design} and
culture of \textbf{Small \& Frequent Releases}.

\textbf{Original post:}
\href{https://www.facebook.com/groups/awsstudygroupfcj/permalink/2229176364513990/?rdid=aT7nNGRCmytCuj4K\#}{Facebook
AWS Study Group FCJ}
